%%%%%%%%%%%%%%%%%%%%%%%%%%%%%%%%%%%%%%%%%%%%%%%%%%%%%%%%%%%%
\section{Quellcode-Listings}%
\index{Listing!Gestaltungsstil}%
\label{sec:Listings}
%%%%%%%%%%%%%%%%%%%%%%%%%%%%%%%%%%%%%%%%%%%%%%%%%%%%%%%%%%%%
%
Zum Einbinden und Formatieren von \index{Code|see{Listing}}Quellcode"=Beispielen
-- sog. \index{Listing}Listings -- wird das Paket \pkg{listings}
\parencite{Hoffmann2014} verwendet.
Das Hervorheben von \index{Schlusselwort@Schlüsselwort}Schlüsselwörtern
wird von LaTeX automatisch erledigt,
wenn die korrekte Sprache des Listings angegeben ist.
Dies geschieht mit Hilfe der Option \printkeyword{language}
oder durch die Angabe eines entsprechend für die jeweilige Sprache
vordefinierten Gestaltungsstils.

Im Befehl \lstinline|\lstset{...}|,
welcher in der Datei \printfilepathquoted{./preambel/preambel.tex} zu finden ist,
kann man einen globalen Stil für alle Listings vorgeben
(welcher jedoch bei Bedarf im Einzelfall überrufen werden kann).
Aktuell ist der etwas weiter oben im Code mit dem Befehl
\lstinline|\lstdefinestyle{...}| vordefinierte Stil
\printkeyword{numbersframed} als Standardgestaltungsstil ausgewählt.
Möchte man die Zeilennummern nicht anzeigen, wäre stattdessen
\printkeyword{nonumbersframed} zu wählen.

%Neben diesen beiden Gestaltungsstilen sind in der Datei
%\printfilepathquoted{./preambel/preambel.tex}
%auch noch \printkeyword{python}, \printkeyword{java}, \printkeyword{C++}
%und \printkeyword{latex} als Gestaltungsstile vordefiniert.
%Bei Bedarf lassen sich dort weitere Stile definieren und auswählen.

Zur Vereinfachung wurden für die \og Programmiersprachen
eigene Umgebungen definiert, die mit \lstinline|\begin{<name>}...\end{<name>}|
direkt verwendet werden können.
Ein Beispiel für die Einbindung eines Java"=Quellcodelistings mit Hilfe einer solchen Umgebung ist in \cref{lst:java-listing} gezeigt.

Man beachte, dass anders als bei Abbildungen und Tabellen
die Bezeichnung (\texttt{caption}) und die Referenzmarke (\texttt{label})
nicht als gesonderte Befehle, sondern als optionale Argumente übergeben werden.
Dies liegt daran, dass ein Listing in der Regel keine Fließumgebung ist,
sondern an der Stelle im Text erscheint, an der sie im Code auch steht.
Ferner folgt ein Listing den ganz normalen Seitenumbruchsregeln.
Das heißt, ein überlanger Code wird einfach umgebrochen. 
Um ein Listing zu einem Fließobjekt zu machen, muss das optionale Argument
\texttt{float=<tbp>} angegeben werden.
Die \index{Platzierung}Platzierungsangabe \enquote{\texttt{h}} für \enquote{hier}
ist nicht erlaubt. Denn dies ist das Standardverhalten ohne \texttt{float}.
%
\begin{latex}[caption={Beispiel eines LaTeX-Quellcodes für ein Java-Listing},label={lst:java-listing}]
\begin{java}[caption={Ein Hello-World Beispiel in Java},
             label={lst:hello-world}]
// Klasse HelloWorld
public class HelloWorld {
  public static void main( String[] args ) {
    System.out.println( "HelloWorld" );
  }
}
\end{java}
\end{latex}
%
Dieses Listing bewirkt das folgende Ergebnis:
%
\begin{java}[caption={Ein Hello-World Beispiel in Java},
             label={lst:hello-world}]
// Klasse HelloWorld
public class HelloWorld {
  public static void main( String[] args ) {
    System.out.println( "HelloWorld" );
  }
}
\end{java}

